\newpage
\section{Quantum Backend}

The quantum backend leverages Qiskit to provide both high-performance statevector simulation and (optionally) access to real IBM Quantum hardware. It interprets the continuous fields of the APEIRON substrate as quantum states, enabling operations like superposition, entanglement, and quantum interference. This backend is a cornerstone of APEIRON-Core's commitment to supporting all available computational matter, from classical to quantum.

\subsection{Statevector Simulation}

The backend uses Qiskit's \texttt{StatevectorSimulator} for exact simulation of quantum circuits. Each field is encoded as a quantum circuit with \(n = \min(d, N_{\text{qubits}})\) qubits, where \(d\) is the field dimension and \(N_{\text{qubits}}\) is configured in the backend. The circuit is initialized to a superposition state via Hadamard gates, and small random rotations are applied to introduce variation:

\begin{lstlisting}[language=Python, caption={Encoding a continuous field as a quantum circuit.}]
qr = QuantumRegister(n_qubits, 'q')
circuit = QuantumCircuit(qr)
for i in range(n_qubits):
    circuit.h(i)
for i in range(n_qubits):
    circuit.rz(random.uniform(0, 2*np.pi), i)
\end{lstlisting}

The resulting statevector is sampled to produce a classical NumPy array that can be consumed by other APEIRON layers.

\subsection{Corrected SWAP Test for Interference}

The interference between two fields is computed as the quantum state overlap \(|\langle \psi_A|\psi_B\rangle|^2\). This is measured using the SWAP test, a standard quantum algorithm that estimates overlap without full state tomography. APEIRON-Core implements a corrected version of the SWAP test that correctly interprets the ancilla qubit's measurement outcome:

\begin{equation}
\mathrm{overlap} = 2 \cdot P(\mathrm{ancilla} = |0\rangle) - 1
\end{equation}

In Qiskit's little-endian bit ordering, the ancilla qubit corresponds to the rightmost bit of the measurement bitstring. The backend correctly sums counts where this bit is \texttt{'0'}:

\begin{lstlisting}[language=Python, caption={Corrected SWAP test implementation.}]
p0 = sum(count for bitstring, count in result.items() if bitstring[-1] == '0') / shots_total
overlap = max(0.0, min(1.0, 2 * p0 - 1))
\end{lstlisting}

This correction addresses a subtle bug that caused overestimation of overlap in earlier implementations.

We use Qiskit Aer 0.13.0 with the \texttt{StatevectorSimulator} (method='statevector') for exact overlap computation. The number of shots for the SWAP test is configurable (default 8192) to balance precision and execution time. To verify the correction, we computed the overlap between identical states (expected 1.0) and orthogonal states (expected 0.0) across 500 random state pairs; the mean absolute error was 0.0023 (Section~10.4).

\subsection{Partial Trace for Entangled Updates}

When entanglement is enabled, the \texttt{field\_update} operation entangles the target field with its past and future states using CNOT gates. The combined statevector of \(3n\) qubits is then partially traced to recover the marginal state of each field. To avoid the performance overhead of Qiskit's \texttt{partial\_trace} function (which requires multiple circuit evaluations), the backend performs the trace directly on the statevector using NumPy reshaping:

\begin{lstlisting}[language=Python, caption={Partial trace via NumPy tensor reshaping.}]
sv = np.array(statevector).reshape([2] * total_qubits)
rho_field = np.zeros((2**n_qubits, 2**n_qubits), dtype=complex)
for i in range(2**n_qubits):
    for j in range(2**n_qubits):
        # trace over past and future qubits
        rho_field[i, j] = sv[:, i, j, ...].sum()
\end{lstlisting}

This optimization reduces the cost of an entangled update from three circuit evaluations to a single evaluation plus \(O(4^n)\) post-processing, a significant improvement for moderate qubit counts (\(n \leq 10\)). Beyond \(n=12\) the statevector size exceeds 64\,GB, necessitating the use of IBM Quantum's 127-qubit \texttt{ibm\_kyiv} device, where we rely on hardware-native SWAP tests and measurement-based state preparation.